\documentclass[11pt]{article}

\usepackage[T1]{fontenc}
\usepackage[utf8]{inputenc}
\usepackage{lmodern}
\usepackage[margin=1.05in]{geometry}
\usepackage{amsmath,amssymb,amsthm}
\usepackage{aliascnt}
\usepackage{microtype}
\usepackage{xcolor}
\usepackage{hyperref}
\usepackage[capitalize,noabbrev]{cleveref}

\hypersetup{
  colorlinks=true,
  linkcolor=blue!45!black,
  citecolor=green!35!black,
  urlcolor=blue!55!black,
  pdftitle={Algebraic Exponential Integrals: Transcendence and Linear Relations},
  pdfauthor={Christopher D. Long; Antoine-Auguste Le Blanc}
}

\numberwithin{equation}{section}
\setlength{\emergencystretch}{3em}

\newtheorem{theorem}{Theorem}[section]

\newaliascnt{proposition}{theorem}
\newtheorem{proposition}[proposition]{Proposition}
\aliascntresetthe{proposition}

\newaliascnt{lemma}{theorem}
\newtheorem{lemma}[lemma]{Lemma}
\aliascntresetthe{lemma}

\newaliascnt{corollary}{theorem}
\newtheorem{corollary}[corollary]{Corollary}
\aliascntresetthe{corollary}

\theoremstyle{definition}
\newaliascnt{definition}{theorem}
\newtheorem{definition}[definition]{Definition}
\aliascntresetthe{definition}

\theoremstyle{remark}
\newaliascnt{remark}{theorem}
\newtheorem{remark}[remark]{Remark}
\aliascntresetthe{remark}

\newcommand{\C}{\mathbb C}
\newcommand{\Qbar}{\overline{\mathbb Q}}
\newcommand{\Aone}{\mathbb A^1}
\newcommand{\Span}{\operatorname{span}}
\newcommand{\Spec}{\operatorname{Spec}}
\newcommand{\ord}{\operatorname{ord}}
\newcommand{\supp}{\operatorname{supp}}
\newcommand{\Tr}{\operatorname{Tr}}
\newcommand{\Mcal}{\mathcal M}
\newcommand{\Ecal}{\mathcal E}
\newcommand{\Rcal}{\mathcal R}
\newcommand{\Lcal}{\mathcal L}
\newcommand{\Dmod}{\mathcal D}
\newcommand{\Ocal}{\mathcal O}

\title{Algebraic Exponential Integrals:\\
Transcendence and Linear Relations}
\author{Christopher D. Long \and Antoine-Auguste Le Blanc\thanks{Le Blanc cryptographic author identifier: \texttt{b6048125\allowbreak{}30b3535b\allowbreak{}23d6dbb0\allowbreak{}f5abe32d\allowbreak{}1ed1a9de\allowbreak{}35fb4f7d\allowbreak{}04d62cff\allowbreak{}08ea19c8}}}
\date{August 2026}

\begin{document}
\maketitle

\begin{center}
\small
\texttt{galizur@gmail.com}
\end{center}

\begin{abstract}
Let $q\in\Qbar[x]$ be nonconstant, and let
$\Delta=\sum_\alpha d_\alpha[\alpha]$ be a finite degree-zero zero-cycle on
$\Aone(\Qbar)$ with coefficients in $\Qbar$.  If $\Gamma$ is any finite
piecewise $C^1$ one-chain in $\C$ with $\partial\Gamma=\Delta$, put
\[
  I_{q,\Delta}(z)=\int_\Gamma e^{zq(x)}\,dx.
\]
We prove that every nonzero function $I_{q,\Delta}$ is a purely
transcendental $E$-function: for every $\xi\in\Qbar^\times$,
$I_{q,\Delta}(\xi)$ is transcendental.  Equivalently,
\[
  I_{q,\Delta}(1)\in\Qbar
  \quad\Longleftrightarrow\quad
  I_{q,\Delta}\equiv0.
\]
In particular, $\int_A^B e^{q(x)}dx$ is transcendental for every pair of
distinct algebraic endpoints.  For a fixed polynomial $q$, we also determine all $\Qbar$-linear
relations among finite families of these values: they are exactly the
specializations of constant-coefficient functional relations, and the latter
are characterized by vanishing of explicit inverse-branch densities on a
finite tree in the target of $q$.

The proof combines Beukers' lifting theorem with a homogeneous moment
connection.  A Fourier--Laplace description proves semisimplicity, while an
explicit residue calculation shows that every rational horizontal
endomorphism is constant.  An algebraic value therefore forces an
exponential-polynomial identity.  Constant-coefficient annihilation and a
Cauchy-transform argument then force the original exponential period
function to vanish identically.
\end{abstract}

\medskip
\noindent\textit{2020 Mathematics Subject Classification.}
Primary 11J81; Secondary 14F10, 34M35.

\noindent\textit{Keywords.}
$E$-functions, exponential period functions, Siegel--Shidlovskii theorem,
Fourier--Laplace transform, polynomial moments, Cauchy transforms,
transcendence, linear relations.

\section{Introduction and main results}

Following Fres\'an and Jossen, we use \emph{exponential period function} for
a parameter-dependent integral of the form $\int_\gamma e^{-zf}\omega$ with
algebraic data \cite{FresanJossen}.  The present functions form a
one-dimensional relative family: $X=\Aone$, the regular function is $-q$,
the algebraic differential form is $dx$, and the integration chain has finite
boundary.  This is compatible with the general terminology of exponential
periods \cite{CommelinHabeggerHuber}.  We refer simply to the polynomial $q$
and to its values and critical values.

\begin{definition}[Zero-cycles and exponential period functions]
A finite zero-cycle on $\Aone(\Qbar)$ with coefficients in $\Qbar$ is a
formal sum
\[
  \Delta=\sum_{\alpha\in S}d_\alpha[\alpha],
  \qquad
  S\subset\Qbar\text{ finite},\quad d_\alpha\in\Qbar.
\]
Its degree is $\deg\Delta=\sum_{\alpha\in S}d_\alpha$.  Suppose
$\deg\Delta=0$.  Choose a finite piecewise $C^1$ one-chain $\Gamma$ in
$\C$, with complex coefficients, such that $\partial\Gamma=\Delta$, and set
\begin{equation}\label{eq:I}
  I_{q,\Delta}(z)=\int_\Gamma e^{zq(x)}\,dx.
\end{equation}
For every $z$, the one-form $e^{zq(x)}dx$ has an entire primitive in $x$.
Thus \eqref{eq:I} is independent of the chosen $\Gamma$, and compactness of
$\Gamma$ shows that $I_{q,\Delta}$ is entire in $z$.
\end{definition}

\begin{theorem}[Algebraic-value rigidity]\label{thm:rigidity}
Let $q\in\Qbar[x]$ be nonconstant, and let $\Delta$ be a finite degree-zero
zero-cycle on $\Aone(\Qbar)$ with coefficients in $\Qbar$.  Then
\begin{equation}\label{eq:rigidity}
  I_{q,\Delta}(1)\in\Qbar
  \quad\Longleftrightarrow\quad
  I_{q,\Delta}(z)\equiv0.
\end{equation}
In the equivalent cases, $I_{q,\Delta}(1)=0$.
\end{theorem}

The fixed evaluation point $1$ entails no loss: a nonzero algebraic
parameter is already part of the polynomial in the exponent.

\begin{corollary}[Pure transcendence and algebraic zeros]
\label{cor:purely-transcendental}
If $I_{q,\Delta}$ is not identically zero, then
\[
  I_{q,\Delta}(\xi)\notin\Qbar
  \qquad(\xi\in\Qbar^\times).
\]
In particular, $I_{q,\Delta}$ has no nonzero algebraic zero.  As an
$E$-function, it is purely transcendental in the standard sense.
\end{corollary}

\begin{proof}
For $\xi\in\Qbar^\times$,
\[
  I_{\xi q,\Delta}(z)=I_{q,\Delta}(\xi z).
\]
The right side is nonzero as an entire function if and only if
$I_{q,\Delta}$ is nonzero.  Apply \cref{thm:rigidity} to the polynomial
$\xi q$.
\end{proof}

\begin{corollary}[A single algebraic exponential integral]
\label{cor:single}
Let $q\in\Qbar[x]$ be nonconstant and let $A,B\in\Qbar$ with $A\ne B$.
Then, for every $\xi\in\Qbar^\times$,
\[
  \int_A^B e^{\xi q(x)}\,dx
\]
is transcendental and nonzero.
\end{corollary}

\begin{proof}
For $\Delta=[B]-[A]$, one has $I_{q,\Delta}(0)=B-A\ne0$, so
$I_{q,\Delta}$ is not identically zero.  Apply
\cref{cor:purely-transcendental}.
\end{proof}

For a fixed $E$-function, Adamczewski and Rivoal gave an algorithm that
computes all nonzero algebraic arguments at which its value is algebraic
\cite{AdamczewskiRivoal}; Bostan, Rivoal, and Salvy made the required
minimal-equation computations effective in practice
\cite{BostanRivoalSalvy}.  The point of \cref{thm:rigidity} is different:
it gives a uniform answer for the entire family indexed by $(q,\Delta)$ and
identifies exactly when the function itself vanishes.  Delaygue's
Lindemann--Weierstrass theorem gives linear independence under a separation
condition on the finite singularity sets of the associated $G$-series
\cite{Delaygue}.  The functions attached to a fixed $q$ arise from one
common moment connection and can have overlapping singularity data; the
results below require no such separation and describe the whole relation
space.

There is classical overlap in special cases.  For example,
\[
  \int_0^A e^{cx^n}\,dx
  =A\,{}_1F_1\!\left(\frac1n;1+\frac1n;cA^n\right),
\]
so monomial exponents belong to the hypergeometric $E$-function setting;
see, for example, Delaygue's discussion of hypergeometric values
\cite{Delaygue}.  Our argument treats arbitrary polynomials, arbitrary
algebraic degree-zero zero-cycles, and all functional relations among a
fixed-$q$ family.  The no-algebraic-zero consequence also connects with
recent work of Fischler and Rivoal on zeros of $E$-functions
\cite{FischlerRivoalZeros}.

The proof has three parts.  First, moments of \eqref{eq:I} satisfy one
explicit inhomogeneous $E$-function system.  Second, the homogeneous
connection is semisimple and every rational horizontal endomorphism is
constant; this converts a polynomial-coefficient relation supplied by
Beukers' theorem into a constant-coordinate relation.  Third, an
exponential-polynomial identity produces vanishing polynomial moments, and
a Cauchy transform on a finite tree forces all inverse-branch densities to
vanish.

\paragraph{Scope of the differential form.}
The fixed form $dx$ is essential.  For $q(x)=x(x-1)$ and
$r(x)=2x^2-x+1$, one has
\[
  r(x)e^{q(x)}=\frac{d}{dx}\bigl(xe^{q(x)}\bigr),
  \qquad \int_0^1 r(x)e^{q(x)}\,dx=1.
\]
Yet the function $\int_0^1 r(x)e^{zq(x)}\,dx$ is not identically zero:
its value at $z=0$ is $7/6$.  Thus \cref{thm:rigidity} does not extend
unchanged to arbitrary polynomial amplitudes.

\section{The moment system}

We use the standard definition: an entire function
$F(z)=\sum_{m\ge0}a_m z^m/m!$ is an $E$-function if its coefficients are
algebraic, $h(a_0,\ldots,a_m)=O(m)$, and $F$ satisfies a nonzero linear
differential equation over $\Qbar[z]$.

Assume throughout this section that $n=\deg q\ge2$, and put
\[
  V=\C[x]_{<n-1}.
\]
For $0\le j\le n-2$, divide
\begin{equation}\label{eq:division}
  q(x)x^j=h_j(x)q'(x)+r_j(x),
  \qquad \deg r_j<n-1.
\end{equation}
Define endomorphisms $\mathsf C,\mathsf D$ of $V$ by
\[
  \mathsf C(x^j)=r_j,
  \qquad
  \mathsf D(x^j)=-h_j'.
\]
Let $C=\mathsf C^{\mathsf T}$ and $D=\mathsf D^{\mathsf T}$ in the
monomial basis $1,x,\ldots,x^{n-2}$.

\begin{lemma}[Residue spectrum]\label{lem:D-spectrum}
The matrix $D$ is triangular, with
\begin{equation}\label{eq:D-spectrum}
  \Spec(D)=
  \left\{-\frac1n,-\frac2n,\ldots,-\frac{n-1}{n}\right\}.
\end{equation}
\end{lemma}

\begin{proof}
Comparison of the leading terms in \eqref{eq:division} gives
\[
  h_j(x)=\frac{x^{j+1}}n+O(x^j).
\]
Consequently,
$-h_j'(x)=-(j+1)x^j/n+O(x^{j-1})$.
\end{proof}

Fix a chain $\Gamma$ with boundary $\Delta$, and put
\begin{equation}\label{eq:U}
  U_j(z)=\int_\Gamma x^j e^{zq(x)}\,dx,
  \qquad
  \mathbf U=(U_0,\ldots,U_{n-2})^{\mathsf T}.
\end{equation}
Let
\begin{equation}\label{eq:Lambda}
  \Lambda=\{0\}\cup q(\supp\Delta),
\end{equation}
with repetitions removed, and define
\[
  \mathbf h(x)=(h_0(x),\ldots,h_{n-2}(x))^{\mathsf T},
  \qquad
  \mathbf b_\lambda=
  \sum_{\substack{\alpha\in\supp\Delta\\q(\alpha)=\lambda}}
     d_\alpha\mathbf h(\alpha)
  \quad(\lambda\in\Lambda).
\]
The empty sum is zero; in particular, $\mathbf b_0=0$ when no endpoint has
$q$-value zero.

\begin{proposition}[Moment system]\label{prop:moment-system}
The moment vector satisfies
\begin{equation}\label{eq:moment-system}
  \mathbf U'
  =\left(C+\frac Dz\right)\mathbf U
   +\frac1z\sum_{\lambda\in\Lambda}
      \mathbf b_\lambda e^{\lambda z}.
\end{equation}
Each $U_j$ is an $E$-function.
\end{proposition}

\begin{proof}
Integration by parts in \eqref{eq:division} gives, for $z\ne0$,
\[
  U_j'
  =\int_\Gamma\left(r_j-\frac{h_j'}z\right)e^{zq}\,dx
   +\frac1z\sum_{\alpha\in\supp\Delta}
      d_\alpha h_j(\alpha)e^{zq(\alpha)},
\]
which is \eqref{eq:moment-system}.  Both sides are meromorphic at $0$, and
the identity holds there by continuation.

Write
\[
  U_j(z)=\sum_{m\ge0}u_{j,m}\frac{z^m}{m!},
  \qquad
  u_{j,m}=\int_\Gamma x^j q(x)^m\,dx,
\]
and
\[
  q(x)^m=\sum_{k=0}^{nm}\gamma_{m,k}x^k.
\]
A polynomial primitive gives the exact formula
\begin{equation}\label{eq:coefficient-formula}
  u_{j,m}=\sum_{k=0}^{nm}\gamma_{m,k}
  \frac{\displaystyle\sum_{\alpha\in\supp\Delta}
        d_\alpha\alpha^{j+k+1}}{j+k+1}.
\end{equation}
Thus every $u_{j,m}$ is algebraic.

We record the arithmetic growth explicitly.  Let $K$ be a number field
containing the coefficients of $q$ and all $\alpha,d_\alpha$.  Uniformly
over embeddings $\sigma:K\hookrightarrow\C$, integration of a primitive
along the segment from $0$ to $\sigma(\alpha)$ gives
\begin{equation}\label{eq:arch-bound}
  |\sigma(u_{j,m})|\le A^{m+j+2}
\end{equation}
for a constant $A\ge1$.  Choose an integer $d\ge1$ such that $d$ times
every coefficient of $q$, every $\alpha$, and every $d_\alpha$ is an
algebraic integer.  For $m\ge0$, set
\begin{equation}\label{eq:common-denominator}
  \delta_m=
  d^{(n+1)m+j+2}\operatorname{lcm}(1,2,\ldots,nm+j+1).
\end{equation}
Formula \eqref{eq:coefficient-formula} shows that
$\delta_m u_{j,\ell}$ is an algebraic integer for every $0\le\ell\le m$:
indeed, $d^\ell\gamma_{\ell,k}$, $d d_\alpha$, and
$d^{j+k+1}\alpha^{j+k+1}$ are integral, while
$k\le n\ell\le nm$, and the total exponent of $d$ needed for each term
satisfies $\ell+j+k+2\le(n+1)m+j+2$.  Since
$\log\operatorname{lcm}(1,\ldots,N)=O(N)$, \eqref{eq:arch-bound} and
\eqref{eq:common-denominator} imply
\[
  h(u_{j,0},\ldots,u_{j,m})=O(m).
\]

Finally, \eqref{eq:moment-system} shows that the $\Qbar(z)$-span of the
$U_j$ and the finitely many functions $e^{\lambda z}$ is stable under
differentiation.  Hence each $U_j$ satisfies a nonzero scalar differential
equation over $\Qbar[z]$.  This proves all defining properties of an
$E$-function.
\end{proof}

When $q(x)=ax+b$ is linear, the explicit identity
\begin{equation}\label{eq:linear-formula}
  I_{q,\Delta}(z)
  =\frac1{az}\sum_{\alpha\in\supp\Delta}
       d_\alpha e^{zq(\alpha)}
  \qquad(z\ne0)
\end{equation}
shows directly that $I_{q,\Delta}$ is also an $E$-function.

We use the following linear form of Beukers' refinement of the
Siegel--Shidlovskii theorem
\cite[Theorem~3.2 and Lemma~3.1]{Beukers}.  No linear independence
hypothesis on the coordinates is required.

\begin{theorem}[Beukers, linear lifting]\label{thm:Beukers}
Let $\mathbf F=(F_1,\ldots,F_N)^{\mathsf T}$ be a vector of $E$-functions
satisfying $\mathbf F'=A(z)\mathbf F$ with
$A(z)\in M_N(\Qbar(z))$, and suppose that $A$ has no pole at $z=1$.
For every relation
\[
  a_1F_1(1)+\cdots+a_NF_N(1)=0,
  \qquad a_i\in\Qbar,
\]
there are polynomials $p_i\in\Qbar[z]$ such that
\[
  \sum_{i=1}^N p_i(z)F_i(z)=0,
  \qquad p_i(1)=a_i\quad(1\le i\le N).
\]
\end{theorem}

\section{The homogeneous moment connection}

The homogeneous part of \eqref{eq:moment-system} is
\begin{equation}\label{eq:homogeneous}
  \mathbf Y'=\left(C+\frac Dz\right)\mathbf Y.
\end{equation}

\begin{proposition}[Semisimplicity]\label{prop:semisimplicity}
The differential module over $\C(z)$ defined by
\eqref{eq:homogeneous} is semisimple.
\end{proposition}

\begin{proof}
Let $\Sigma$ be the finite set of critical values of $q$, and put
$V_0=\Aone_t\setminus\Sigma$.  Over $V_0$, the finite map
$q:q^{-1}(V_0)\to V_0$ defines a rank-$n$ permutation local system
$\Lcal$.  Its monodromy group is finite, so $\Lcal$ is semisimple by
Maschke's theorem.  The trace map splits it as
\[
  \Lcal\simeq\C_{V_0}\oplus\Lcal^0,
  \qquad \Lcal^0=\ker(\Tr:\Lcal\to\C_{V_0}).
\]

The map $q$ is finite, hence proper and small.  The small-map
intermediate-extension identity
\cite[Remark~4.2.4]{deCataldoMigliorini} gives
\begin{equation}\label{eq:small-map}
  Rq_*\C_{\Aone_x}[1]\simeq\operatorname{IC}(\Lcal).
\end{equation}
Intermediate extension preserves direct sums and takes irreducible local
systems to simple perverse sheaves.  Thus the semisimplicity of $\Lcal$
makes this perverse sheaf semisimple.  Under the Riemann--Hilbert
correspondence and
its compatibility with proper direct image, the regular holonomic
$\Dmod$-module $q_+\Ocal_{\Aone_x}$ is therefore semisimple; the constant
summand of $\Lcal$ corresponds to $\Ocal_{\Aone_t}$
\cite{HTT}.

Use a Fourier variable $\tau$ and the Weyl-algebra convention
\begin{equation}\label{eq:Fourier-convention}
  t\longmapsto-\partial_\tau,
  \qquad \partial_t\longmapsto\tau.
\end{equation}
This is the inverse convention to the transform with kernel $e^{t\tau}$
in \cite[\S1.3.b]{SabbahFourier}; our convention corresponds to the
kernel $e^{-t\tau}$.
Fourier transformation is an exact equivalence on holonomic Weyl modules
and preserves simple objects and direct sums.  The transform of
$\Ocal_{\Aone_t}$ is supported at $\tau=0$, hence disappears after
localization to $\C(\tau)$.  A nonzero generic localization of a simple
Weyl module remains irreducible.  Indeed, its
$\C[\tau]\setminus\{0\}$-torsion is a Weyl submodule: if
$p(\tau)m=0$, then $p(\tau)^2\partial_\tau m=0$.  Thus a simple module
whose localization is nonzero embeds in that localization.  Any nonzero
$\C(\tau)$-differential submodule meets the original module after clearing
a denominator, and simplicity then forces equality.  It follows that the
generic Fourier transform of $q_+\Ocal_{\Aone_x}$ is semisimple.

For completeness, we identify this generic transform explicitly.  Factor
$q$ as its graph embedding followed by projection to the $t$-line.  A graph
generator $\delta_q$ satisfies
\[
  (t-q(x))\delta_q=0,
  \qquad
  (\partial_x+q'(x)\partial_t)\delta_q=0.
\]
After \eqref{eq:Fourier-convention}, the transformed generator satisfies
\[
  (\partial_\tau+q(x))\widehat\delta_q=0,
  \qquad
  (\partial_x+\tau q'(x))\widehat\delta_q=0.
\]
In particular,
$\partial_x(p\widehat\delta_q)=(p'-\tau q'p)\widehat\delta_q$:
the relation on the generator induces a minus sign in the differential
on coefficient polynomials.  The generic direct image along the $x$-line
is therefore the degree-zero cohomology of
\[
  \C(\tau)[x]
  \xrightarrow{\ \partial_x-\tau q'(x)\ }
  \C(\tau)[x],
\]
with $\tau$-connection induced by $\partial_\tau-q(x)$.  The displayed
map is injective: for nonzero $p$, the term $-\tau q'p$ has larger
$x$-degree than $p'$.  Hence the generic transform is
\begin{equation}\label{eq:generic-Fourier}
  \C(\tau)[x]\big/
  (\partial_x-\tau q'(x))\C(\tau)[x],
  \qquad
  \nabla_{\partial_\tau}[p]=[\partial_\tau p-qp].
\end{equation}
This is the algebraic Fourier--Laplace direct-image construction of
\cite[\S1.3.b]{SabbahFourier}, with the inverse convention fixed in
\eqref{eq:Fourier-convention}.

Pulling \eqref{eq:generic-Fourier} back by $\tau=-z$ gives
\begin{equation}\label{eq:twisted-module}
  \Mcal_q=
  \C(z)[x]\big/(\partial_x+zq'(x))\C(z)[x],
  \qquad
  \nabla_{\partial_z}[p]=[\partial_zp+qp].
\end{equation}
The connection descends to this quotient because
\begin{equation}\label{eq:twisted-commutator}
  [\partial_z+q,\partial_x+zq']=q'-q'=0.
\end{equation}
The classes of $1,x,\ldots,x^{n-2}$ form a basis: the relation
$[zq'p]=-[p']$ reduces every higher monomial, while degree comparison proves
independence.  In this basis, \eqref{eq:division} gives the connection
matrix $\mathsf C+\mathsf D/z$.  The period system
\eqref{eq:homogeneous} is the corresponding dual system, and duality
preserves semisimplicity.
\end{proof}

\begin{proposition}[Rational endomorphism rigidity]
\label{prop:endomorphism-rigidity}
If $P(z)\in M_{n-1}(\C(z))$ satisfies
\begin{equation}\label{eq:endomorphism}
  P'=\left[C+\frac Dz,P\right],
\end{equation}
then $P$ is constant.  The horizontal endomorphism algebra of
\eqref{eq:homogeneous} is therefore
\[
  Z(C)\cap Z(D).
\]
The same constancy conclusion holds for the dual module.
\end{proposition}

\begin{proof}
A pole at $z_0\ne0$ is impossible: differentiation raises its order by
one, whereas the right side of \eqref{eq:endomorphism} does not.  At $0$,
a leading term $z^{-m}P_{-m}$ with $m\ge1$ would give
\[
  -mP_{-m}=[D,P_{-m}].
\]
By \eqref{eq:D-spectrum}, every eigenvalue of
$\operatorname{ad}D$ lies strictly between $-1$ and $1$, so the negative
integer $-m$ cannot occur.  Thus $P$ has no finite pole and is polynomial.

It remains to exclude positive degree.  Transpose
\eqref{eq:endomorphism} and write $Q=P^{\mathsf T}$.  The operator
$\mathsf C$ is multiplication by $q$ on $\C[x]/(q')$.  For a critical
point $\tau$, put
\[
  m_\tau=\ord_\tau q',
  \qquad \lambda_\tau=q(\tau),
\]
and define
\[
  g_{\tau,k}(x)=\frac{q'(x)}{(x-\tau)^k},
  \qquad 1\le k\le m_\tau.
\]
The Chinese remainder theorem shows that all the $g_{\tau,k}$ form a
basis of $V$.  Moreover,
\[
  \mathsf Cg_{\tau,k}=\lambda_\tau g_{\tau,k},
\]
because
\[
  qg_{\tau,k}
  =\frac{q-\lambda_\tau}{(x-\tau)^k}q'
   +\lambda_\tau g_{\tau,k}.
\]
Thus $\mathsf C$ is semisimple, with eigenspaces indexed by critical
values.

Let $E_\lambda$ be the $\lambda$-eigenspace of $\mathsf C$, let
$\pi_\lambda$ be its spectral projection, and put
\[
  \overline D_\lambda
  =\pi_\lambda\mathsf D|_{E_\lambda}.
\]
Fix $\tau$ with $\lambda_\tau=\lambda$, and write $m=m_\tau$.
Define
\[
  F_{\tau,0}=0,\qquad
  F_{\tau,k}=\Span_\C\{g_{\tau,1},\ldots,g_{\tau,k}\}
  \quad(1\le k\le m).
\]
We claim the explicit flag relation
\begin{equation}\label{eq:local-flag}
  \overline D_\lambda g_{\tau,k}
  +\frac{m+1-k}{m+1}g_{\tau,k}
  \in F_{\tau,k-1}
  \qquad(1\le k\le m).
\end{equation}
To prove it, write $y=x-\tau$ and
$q-\lambda=ay^{m+1}+O(y^{m+2})$, with $a\ne0$.  Euclidean division gives
\[
  \mathsf Dg_{\tau,k}
  =-\left(\frac{q-\lambda}{(x-\tau)^k}\right)',
\]
so
\begin{align*}
  g_{\tau,k}&=(m+1)a y^{m-k}+O(y^{m-k+1}),\\
  \mathsf Dg_{\tau,k}&=-(m+1-k)a y^{m-k}+O(y^{m-k+1}).
\end{align*}
Their combination in \eqref{eq:local-flag} therefore vanishes at $\tau$
to order at least $m-k+1$.  If $\sigma\ne\tau$ has the same critical
value, then $(x-\tau)^{-k}$ is regular at $\sigma$, while $q-\lambda$
vanishes there to order $m_\sigma+1$.  Hence
$\mathsf Dg_{\tau,k}$ vanishes at $\sigma$ to order at least $m_\sigma$,
and its component in that factor of $\C[x]/(q')$ is zero.

Under the Chinese remainder decomposition, $F_{\tau,k-1}$ consists
precisely of the classes supported in the $\tau$-factor and divisible
there by $y^{m-k+1}$.  Projection by $\pi_\lambda$ removes all factors
over other critical values.  The preceding vanishing statements thus
prove \eqref{eq:local-flag}.  In particular, every $F_{\tau,k}$ is
$\overline D_\lambda$-invariant.  On each $F_{\tau,m}$ the matrix in the
ordered basis $g_{\tau,1},\ldots,g_{\tau,m}$ is triangular, with diagonal
\begin{equation}\label{eq:local-spectrum}
  -\frac{m}{m+1},
  -\frac{m-1}{m+1},\ldots,-\frac1{m+1}.
\end{equation}
As $\tau$ ranges over the critical points with value $\lambda$, these
lists give the full spectrum of $\overline D_\lambda$, with
multiplicities.  In particular, all its eigenvalues lie in $(-1,0)$,
including when distinct critical points have the same value and different
ramification indices.

Suppose
\[
  Q(z)=Q_dz^d+Q_{d-1}z^{d-1}+\cdots,
  \qquad d\ge1,\quad Q_d\ne0.
\]
The top two coefficients of the transposed equation give
\begin{equation}\label{eq:top-coefficients}
  [Q_d,\mathsf C]=0,
  \qquad
  dQ_d=[Q_{d-1},\mathsf C]+[Q_d,\mathsf D].
\end{equation}
The first identity makes $Q_d$ block diagonal for the eigenspace
decomposition of $\mathsf C$.  Choose a critical value $\lambda$ for
which $Q_{d,\lambda}\ne0$; the value $\lambda=0$ is allowed.
Projecting the second identity to this block yields
\[
  dQ_{d,\lambda}
  =[Q_{d,\lambda},\overline D_\lambda].
\]
After triangularizing $\overline D_\lambda$, the eigenvalues of the
commutator operator on the right are pairwise differences of the numbers in
\eqref{eq:local-spectrum}.  They all lie strictly between $-1$ and $1$,
so none equals the positive integer $d$.  This contradiction proves that
$d=0$.

A constant solution of \eqref{eq:endomorphism} commutes with both $C$ and
$D$, and the converse is immediate.  The calculation for a right-acting
endomorphism of the dual gives the same equation.
\end{proof}

\begin{corollary}[Constant splitting]\label{cor:constant-splitting}
Every differential submodule of the dual of \eqref{eq:homogeneous} has the
form
\[
  W\otimes_\C\C(z)
\]
for a constant subspace $W\subseteq\C^{1\times(n-1)}$ invariant under
right multiplication by $C$ and $D$.
\end{corollary}

\begin{proof}
By \cref{prop:semisimplicity}, a differential submodule has a differential
complement and hence a horizontal idempotent projector.  By
\cref{prop:endomorphism-rigidity}, this projector is constant.  Its image
is the required subspace $W$; substituting the constant projector into
\eqref{eq:endomorphism} shows that $W$ is $C$- and $D$-invariant.
\end{proof}

\section{From an algebraic value to an exponential polynomial}

Put
\[
  \Ecal=\Span_{\C(z)}\{e^{\lambda z}:\lambda\in\Lambda\}.
\]

\begin{proposition}\label{prop:exp-poly}
If $U_0(1)\in\Qbar$, then
\begin{equation}\label{eq:exp-poly}
  U_0(z)=\sum_{\lambda\in\Lambda}p_\lambda(z)e^{\lambda z}
\end{equation}
with $p_\lambda\in\Qbar[z]$.
\end{proposition}

\begin{proof}
Adjoin the exponential coordinates $e^{\lambda z}$ to the moment vector.
By \eqref{eq:moment-system}, the resulting vector of $E$-functions
satisfies a homogeneous block system whose only finite pole is at $z=0$.
Let $a=U_0(1)\in\Qbar$.  Applying \cref{thm:Beukers} to
$U_0(1)-a\cdot1=0$ gives a polynomial row
$\boldsymbol\phi(z)$ and polynomials $c_\lambda(z)$ such that
\begin{equation}\label{eq:lifted}
  \boldsymbol\phi(z)\mathbf U(z)
  +\sum_{\lambda\in\Lambda}c_\lambda(z)e^{\lambda z}=0,
  \qquad
  \boldsymbol\phi(1)=(1,0,\ldots,0).
\end{equation}

Define
\[
  \Rcal=
  \{\boldsymbol\psi\in\C(z)^{1\times(n-1)}:
    \boldsymbol\psi\mathbf U\in\Ecal\}.
\]
If $A(z)=C+D/z$ and
$\mathbf f=z^{-1}\sum_\lambda\mathbf b_\lambda e^{\lambda z}$, then
\[
  (\boldsymbol\psi\mathbf U)'
  =(\boldsymbol\psi'+\boldsymbol\psi A)\mathbf U
   +\boldsymbol\psi\mathbf f.
\]
Since $\Ecal$ is stable under differentiation, this proves that $\Rcal$
is a differential submodule of the dual homogeneous connection.  By
\cref{cor:constant-splitting},
$\Rcal=W\otimes_\C\C(z)$ for a constant $C$- and $D$-invariant subspace
$W$.  Equation \eqref{eq:lifted} puts $\boldsymbol\phi$ in $\Rcal$.
Every constant linear form annihilating $W$ therefore annihilates
$\boldsymbol\phi(z)$ identically, and hence also annihilates
$\boldsymbol\phi(1)$.  Thus $\boldsymbol\phi(1)\in W$, without
specializing any rational basis.  Hence the constant row
$\mathbf e_0=(1,0,\ldots,0)$ belongs to $W$, and
$U_0=\mathbf e_0\mathbf U$ belongs to $\Ecal$.

Choose a full-row-rank constant matrix $S$ whose rows form a basis of $W$,
and put $\mathbf V=S\mathbf U$.  There are constant matrices $C_W,D_W$
such that
\[
  SC=C_WS,
  \qquad SD=D_WS.
\]
Every component of $\mathbf V$ belongs to $\Ecal$.  The functions
$e^{\lambda z}$ are linearly independent over $\C(z)$, so there is a unique
expansion
\[
  \mathbf V=\sum_{\lambda\in\Lambda}
  \mathbf r_\lambda(z)e^{\lambda z},
  \qquad \mathbf r_\lambda\in\C(z)^{\dim W}.
\]
Equating exponential coefficients in the restricted moment system gives
\begin{equation}\label{eq:r-system}
  \mathbf r_\lambda'
  =\left(C_W-\lambda I+\frac{D_W}{z}\right)
     \mathbf r_\lambda
   +\frac{S\mathbf b_\lambda}{z}.
\end{equation}
A pole away from $0$ is impossible by order comparison.  If
$\mathbf r_\lambda=z^{-m}v+O(z^{-m+1})$ at $0$, with $m\ge1$, then
\eqref{eq:r-system} gives
\[
  -mv=D_Wv.
\]
The eigenvalues of the restriction $D_W$ are among those of $D$, hence lie
strictly between $-1$ and $0$ by \eqref{eq:D-spectrum}; none is a negative
integer.  Thus every $\mathbf r_\lambda$ is polynomial, and so are the
coefficients $p_\lambda$ in \eqref{eq:exp-poly}.

It remains to show that their coefficients are algebraic.  Choose an
integer $d\ge0$ bounding the degrees of all the $p_\lambda$, write
$p_\lambda(z)=\sum_{k=0}^d a_{\lambda,k}z^k$, and put
$M=(d+1)|\Lambda|$.  Differentiating \eqref{eq:exp-poly} at $0$ gives
\begin{equation}\label{eq:finite-descent}
  U_0^{(m)}(0)
  =\sum_{\lambda\in\Lambda}\sum_{k=0}^d
     a_{\lambda,k}
     \left.\frac{d^k}{dt^k}t^m\right|_{t=\lambda}
  \qquad(0\le m<M).
\end{equation}
The square coefficient matrix has algebraic entries and is invertible.
Indeed, its transpose is the Hermite evaluation map
\[
  \Qbar[t]_{<M}\longrightarrow\Qbar^M,
  \qquad f\longmapsto
  \bigl(f^{(k)}(\lambda)\bigr)_{\lambda\in\Lambda,\,0\le k\le d}.
\]
A polynomial in its kernel is divisible by
$\prod_{\lambda\in\Lambda}(t-\lambda)^{d+1}$, of degree $M$, so must
vanish.  Thus the map is injective and hence invertible.  All derivatives
$U_0^{(m)}(0)$ are algebraic by \eqref{eq:coefficient-formula}; solving
\eqref{eq:finite-descent} yields $a_{\lambda,k}\in\Qbar$.  Therefore
$p_\lambda\in\Qbar[z]$.
\end{proof}

\section{Moments, Cauchy transforms, and functional relations}

We first isolate the geometric representation used by the analytic
argument.

\begin{lemma}[Tree representation]\label{lem:tree-representation}
Let $q\in\C[x]$ be nonconstant and let $\Delta$ be a finite degree-zero
zero-cycle on $\C$ with complex coefficients.  There exist a finite
embedded tree $T$ in the $t$-plane, containing all finite critical values
of $q$ and all values $q(\alpha)$ for $\alpha\in\supp\Delta$, and a
one-chain $\widetilde\Gamma$ in the graph $q^{-1}(T)$ such that
$\partial\widetilde\Gamma=\Delta$.

Orient every open edge $e$ of $T$.  Over $e$, write
$\widetilde\Gamma$ as a linear combination of analytic inverse branches
$x_{e,i}$ of $q$, with coefficients $f_{e,i}$, and set
\begin{equation}\label{eq:edge-density}
  \psi_e(t)=\sum_i f_{e,i}x_{e,i}'(t)
  =\sum_i\frac{f_{e,i}}{q'(x_{e,i}(t))}.
\end{equation}
Then $\psi_e$ is analytic on the open edge, has at worst integrable
singularities at its endpoints, and for every entire function $F$,
\begin{equation}\label{eq:edge-decomposition}
  \int_\Gamma F(q(x))\,dx
  =\sum_e\int_e F(t)\psi_e(t)\,dt,
\end{equation}
where $\Gamma$ is any one-chain with boundary $\Delta$.
\end{lemma}

\begin{proof}
Choose a regular value $c$ of $q$, distinct from all endpoint values, and
join $c$ to the distinct finite critical values and endpoint values by
simple analytic arcs with pairwise disjoint interiors.  Their union is a
finite tree $T$.  The inverse image of the critical-value arms is the
standard polynomial constellation.  It is connected because the monodromy
of the polynomial covering on a regular fiber is transitive
\cite[\S2.1]{PakovichMuzychuk}.  The remaining arms attach to this connected
subgraph, so $q^{-1}(T)$ is connected.

Subdivide the graph so that every point of $\supp\Delta$ is a vertex.  Fix
a vertex $x_*$ and choose an edge path $P_\alpha$ from $x_*$ to each
$\alpha\in\supp\Delta$.  Then
\[
  \widetilde\Gamma=
  \sum_{\alpha\in\supp\Delta}d_\alpha P_\alpha
\]
has boundary $\Delta$, because $\deg\Delta=0$.  The inverse-branch
description on open edges gives \eqref{eq:edge-density}.  Near a ramified
endpoint of index $r$, one has
$x_{e,i}'(t)=O(|t-t_0|^{1/r-1})$, which is integrable.

The one-form $F(q(x))dx$ is entire in $x$, so its integral depends only on
the boundary.  Replacing $\Gamma$ by $\widetilde\Gamma$ and changing
variables $t=q(x)$ edge by edge gives \eqref{eq:edge-decomposition}.
\end{proof}

\begin{theorem}[Moment rigidity and the edge-density criterion]
\label{thm:moment-rigidity}
In the setting of \cref{lem:tree-representation}, let
$0\ne H\in\C[t]$.  If
\begin{equation}\label{eq:moments}
  \int_\Gamma q(x)^mH(q(x))\,dx=0
  \qquad(m\ge0),
\end{equation}
then every density $\psi_e$ in \eqref{eq:edge-density} vanishes
identically.  Consequently,
\[
  I_{q,\Delta}(z)\equiv0.
\]
For any fixed tree representation,
\begin{equation}\label{eq:edge-criterion}
  I_{q,\Delta}\equiv0
  \quad\Longleftrightarrow\quad
  \psi_e\equiv0\text{ for every open edge }e\subset T.
\end{equation}
\end{theorem}

\begin{proof}
Every differential $q(x)^mH(q(x))dx$ is polynomial and has a polynomial
primitive, so \eqref{eq:moments} remains true after replacing $\Gamma$ by
$\widetilde\Gamma$.  For $s\in\C\setminus T$, define
\[
  \mathcal C_H(s)=
  \int_{\widetilde\Gamma}
  \frac{H(q(x))}{s-q(x)}\,dx.
\]
For sufficiently large $|s|$, the geometric series for
$(s-q(x))^{-1}$ converges uniformly on the compact chain, and
\eqref{eq:moments} gives $\mathcal C_H(s)=0$.  The complement of a finite
embedded tree is connected, so the identity theorem yields
\begin{equation}\label{eq:Cauchy-zero}
  \mathcal C_H\equiv0\qquad\text{on }\C\setminus T.
\end{equation}
By \eqref{eq:edge-decomposition},
\begin{equation}\label{eq:edge-Cauchy}
  \mathcal C_H(s)
  =\sum_e\int_e\frac{H(t)\psi_e(t)}{s-t}\,dt.
\end{equation}

Fix an interior point $t_0$ of an edge $e$ with $H(t_0)\ne0$, and choose
a compact analytic subarc $\eta\Subset e$ containing $t_0$.  Contributions
to \eqref{eq:edge-Cauchy} from outside $\eta$ extend holomorphically across
a small disk about $t_0$.  Put $\rho(t)=H(t)\psi_e(t)$ and
\[
  J(s)=\int_\eta\frac{\rho(t)}{s-t}\,dt.
\]
If $a,b$ are the oriented endpoints of $\eta$, then locally
\[
  J(s)=
  \rho(s)\bigl(\log(s-a)-\log(s-b)\bigr)
  -\int_\eta\frac{\rho(t)-\rho(s)}{t-s}\,dt.
\]
The second term is holomorphic across $\eta$, while the two lateral
branches of the logarithmic term differ by
$2\pi i\varepsilon_e\rho(t_0)$, where
$\varepsilon_e\in\{\pm1\}$.  By \eqref{eq:Cauchy-zero}, both lateral
values of the full Cauchy transform are zero.  Hence
$H(t_0)\psi_e(t_0)=0$, and therefore $\psi_e(t_0)=0$.  The point $t_0$ was
arbitrary outside the finite zero set of $H$; analyticity gives
$\psi_e\equiv0$ on every open edge.

Taking $F(t)=e^{zt}$ in \eqref{eq:edge-decomposition} now gives
$I_{q,\Delta}\equiv0$.  Conversely, if $I_{q,\Delta}\equiv0$, its Taylor
coefficients give \eqref{eq:moments} with $H=1$, so the preceding argument
forces all densities to vanish.  The reverse implication in
\eqref{eq:edge-criterion} follows directly from
\eqref{eq:edge-decomposition}.
\end{proof}

This argument is a local Cauchy-transform form of techniques used in the
polynomial moment problem and for Cauchy-type integrals of algebraic
functions \cite{PakovichMuzychuk,PakovichRoytvarfYomdin,GavrilovPakovich}.

\section{Proofs and linear relations}

\begin{proof}[Proof of \cref{thm:rigidity}]
The reverse implication is immediate.  Assume that
$I_{q,\Delta}(1)=\gamma\in\Qbar$.

Suppose first that $q(x)=ax+b$ is linear.  Formula
\eqref{eq:linear-formula} at $z=1$ gives
\begin{equation}\label{eq:linear-LW}
  \sum_{\alpha\in\supp\Delta}d_\alpha e^{q(\alpha)}
  -a\gamma e^0=0.
\end{equation}
The values $q(\alpha)$ are pairwise distinct.  Combine the two coefficients
of $e^0$ if one endpoint satisfies $q(\alpha)=0$.  The linear form of the
Lindemann--Weierstrass theorem \cite[Chapter~1]{Baker} then makes every
coefficient in \eqref{eq:linear-LW} zero.  Thus
$d_\alpha=0$ for all endpoints with $q(\alpha)\ne0$.  If
$q(\alpha_0)=0$, then $d_{\alpha_0}=a\gamma$, while
$\deg\Delta=0$ gives $d_{\alpha_0}=0$ and hence $\gamma=0$.  If no such
endpoint exists, the coefficient of $e^0$ gives $\gamma=0$ directly.
Therefore $\Delta=0$ and $I_{q,\Delta}\equiv0$.

Now suppose $\deg q\ge2$.  By \cref{prop:exp-poly},
\[
  U_0(z)=I_{q,\Delta}(z)
  =\sum_{\lambda\in\Lambda}p_\lambda(z)e^{\lambda z}.
\]
If $U_0\equiv0$, there is nothing to prove.  Otherwise set
\[
  N=1+\max_{p_\lambda\ne0}\deg p_\lambda,
  \qquad
  L=\prod_{\lambda\in\Lambda}(\partial_z-\lambda)^N,
\]
and
\[
  H(t)=\prod_{\lambda\in\Lambda}(t-\lambda)^N.
\]
Then $LU_0=0$, and differentiation under the integral sign gives
\[
  0=LU_0(z)
  =\int_\Gamma H(q(x))e^{zq(x)}\,dx.
\]
Comparing Taylor coefficients at $z=0$ gives \eqref{eq:moments}.  Apply
\cref{thm:moment-rigidity}.
\end{proof}

\begin{corollary}[Transfer of linear relations]\label{cor:transfer}
Let $q\in\Qbar[x]$ be nonconstant, let
$\Delta_1,\ldots,\Delta_r$ be finite degree-zero zero-cycles on
$\Aone(\Qbar)$ with coefficients in $\Qbar$, and put
$I_j(z)=I_{q,\Delta_j}(z)$.  For $a_0,\ldots,a_r\in\Qbar$,
\begin{equation}\label{eq:transfer}
  a_0+\sum_{j=1}^r a_jI_j(1)=0
\end{equation}
if and only if
\[
  a_0=0,
  \qquad
  \sum_{j=1}^r a_jI_j(z)\equiv0.
\]
Consequently,
\[
 \dim_{\Qbar}\Span_{\Qbar}\{1,I_1(1),\ldots,I_r(1)\}
 =1+\dim_{\Qbar}\Span_{\Qbar}\{I_1,\ldots,I_r\}.
\]
\end{corollary}

\begin{proof}
Put $\Delta=\sum_j a_j\Delta_j$.  Relation \eqref{eq:transfer} says
$I_{q,\Delta}(1)=-a_0\in\Qbar$.  By \cref{thm:rigidity},
$I_{q,\Delta}\equiv0$, and then its value at $1$ gives $a_0=0$.  The
converse and the dimension identity are immediate.
\end{proof}

The same equivalence holds with the common evaluation point $1$ replaced by
any $\xi\in\Qbar^\times$: replace $q$ by $\xi q$ and then rescale the
function variable.

\begin{corollary}[Geometric description of the relation space]
\label{cor:geometric-relations}
In the setting of \cref{cor:transfer}, choose one tree $T$ containing all
critical values of $q$ and all values
$q(\supp\Delta_j)$.  For each $j$, choose a graph chain representing
$\Delta_j$ as in \cref{lem:tree-representation}, and denote its density on
an open edge $e$ by $\psi_{j,e}$.  Then
\[
  a_0+\sum_{j=1}^r a_jI_j(1)=0
\]
if and only if
\[
  a_0=0,
  \qquad
  \sum_{j=1}^r a_j\psi_{j,e}\equiv0
  \quad\text{for every open edge }e\subset T.
\]
Thus all $\Qbar$-linear relations among the values are determined by a
finite family of inverse-branch identities on $T$.
\end{corollary}

\begin{proof}
The graph chain for $\Delta=\sum_j a_j\Delta_j$ is the same linear
combination of the chosen graph chains, and its density on $e$ is
$\sum_j a_j\psi_{j,e}$.  Combine \cref{cor:transfer} with the exact
criterion \eqref{eq:edge-criterion}.
\end{proof}

\begin{corollary}[Linear independence over distinct regular values]
\label{cor:regular}
Let $q\in\Qbar[x]$ be nonconstant, and let $a_0,\ldots,a_r\in\Qbar$ be
such that $q(a_0),\ldots,q(a_r)$ are pairwise distinct regular values of
$q$.  Then, for every $\xi\in\Qbar^\times$,
\[
  1,\quad
  \int_{a_0}^{a_1}e^{\xi q(x)}\,dx,\quad\ldots,\quad
  \int_{a_0}^{a_r}e^{\xi q(x)}\,dx
\]
are linearly independent over $\Qbar$.
\end{corollary}

\begin{proof}
Absorb $\xi$ into $q$.  It suffices to prove that the functions
\[
  I_j(z)=\int_{a_0}^{a_j}e^{zq(x)}\,dx,
  \qquad 1\le j\le r,
\]
are linearly independent.  A functional relation produces a degree-zero
zero-cycle
\[
  \Delta=\sum_{j=1}^r c_j[a_j]
  -\left(\sum_{j=1}^r c_j\right)[a_0]
\]
whose exponential period function vanishes identically.  Choose the tree
so that each $q(a_i)$ is the endpoint of a separate terminal arm.  Over the
arm ending at $q(a_i)$, regularity and distinctness imply that the only
nonzero terminal boundary coefficient can occur on the lift ending at
$a_i$.  The density there is therefore
\[
  \frac{d_i}{q'(x_i(t))},
\]
where $d_i$ is the coefficient of $[a_i]$ in $\Delta$.  By
\eqref{eq:edge-criterion}, this density vanishes, so $d_i=0$ for every
$i$.  Hence all $c_j=0$.  Apply \cref{cor:transfer}.
\end{proof}

\section*{Acknowledgments and AI-assisted research disclosure}

This manuscript was developed through interactive work between Christopher
D. Long and ChatGPT 5.6 Pro.  ChatGPT 5.6 Pro assisted in proof discovery,
organization, symbolic checking, reference verification, adversarial
auditing, and drafting.  Claude Opus 5 contributed to the
semisimple-projector strategy and independent audits.  The AI systems are
not authors.  The authors bear full responsibility for all statements,
proofs, citations, and any remaining errors.  This draft has not yet
undergone independent peer review or formal verification.

\small
\begin{thebibliography}{99}

\bibitem{AdamczewskiRivoal}
B.~Adamczewski and T.~Rivoal,
\emph{Exceptional values of $E$-functions at algebraic points},
Bull. Lond. Math. Soc. \textbf{50} (2018), no.~4, 697--708.
\href{https://doi.org/10.1112/blms.12168}{doi:10.1112/blms.12168}.

\bibitem{Baker}
A.~Baker,
\emph{Transcendental Number Theory},
Cambridge Mathematical Library,
Cambridge University Press, Cambridge, 1990.

\bibitem{Beukers}
F.~Beukers,
\emph{A refined version of the Siegel--Shidlovskii theorem},
Ann. of Math. (2) \textbf{163} (2006), no.~1, 369--379.
\href{https://doi.org/10.4007/annals.2006.163.369}{doi:10.4007/annals.2006.163.369}.

\bibitem{BostanRivoalSalvy}
A.~Bostan, T.~Rivoal, and B.~Salvy,
\emph{Minimization of differential equations and algebraic values of
$E$-functions},
Math. Comp. \textbf{93} (2024), no.~347, 1427--1472.
\href{https://doi.org/10.1090/mcom/3912}{doi:10.1090/mcom/3912}.

\bibitem{CommelinHabeggerHuber}
J.~Commelin, P.~Habegger, and A.~Huber,
\emph{Exponential periods and o-minimality},
arXiv:2007.08280v4 (2025), to appear in
Ann. Inst. Fourier (Grenoble).

\bibitem{deCataldoMigliorini}
M.~A.~A. de~Cataldo and L.~Migliorini,
\emph{The decomposition theorem, perverse sheaves and the topology of
algebraic maps},
Bull. Amer. Math. Soc. (N.S.) \textbf{46} (2009), no.~4, 535--633.
\href{https://doi.org/10.1090/S0273-0979-09-01260-9}{doi:10.1090/S0273-0979-09-01260-9}.

\bibitem{Delaygue}
\'E.~Delaygue,
\emph{A Lindemann--Weierstrass theorem for $E$-functions},
J. Reine Angew. Math. \textbf{820} (2025), 75--85.
\href{https://doi.org/10.1515/crelle-2024-0090}{doi:10.1515/crelle-2024-0090}.

\bibitem{FischlerRivoalZeros}
S.~Fischler and T.~Rivoal,
\emph{Zeros of $E$-functions and of exponential polynomials defined over
$\overline{\mathbb Q}$},
arXiv:2503.20345 (2025).

\bibitem{FresanJossen}
J.~Fres\'an and P.~Jossen,
\emph{A non-hypergeometric $E$-function},
Ann. of Math. (2) \textbf{194} (2021), no.~3, 903--942.
\href{https://doi.org/10.4007/annals.2021.194.3.7}{doi:10.4007/annals.2021.194.3.7}.

\bibitem{GavrilovPakovich}
L.~Gavrilov and F.~Pakovich,
\emph{Moments on Riemann surfaces and hyperelliptic Abelian integrals},
Comment. Math. Helv. \textbf{89} (2014), no.~1, 125--155.
\href{https://doi.org/10.4171/CMH/314}{doi:10.4171/CMH/314}.

\bibitem{HTT}
R.~Hotta, K.~Takeuchi, and T.~Tanisaki,
\emph{$D$-Modules, Perverse Sheaves, and Representation Theory},
Progress in Mathematics, vol.~236,
Birkh\"auser Boston, Boston, MA, 2008.
\href{https://doi.org/10.1007/978-0-8176-4523-6}{doi:10.1007/978-0-8176-4523-6}.

\bibitem{PakovichMuzychuk}
F.~Pakovich and M.~Muzychuk,
\emph{Solution of the polynomial moment problem},
Proc. Lond. Math. Soc. (3) \textbf{99} (2009), no.~3, 633--657.
\href{https://doi.org/10.1112/plms/pdp010}{doi:10.1112/plms/pdp010}.

\bibitem{PakovichRoytvarfYomdin}
F.~Pakovich, N.~Roytvarf, and Y.~Yomdin,
\emph{Cauchy-type integrals of algebraic functions},
Israel J. Math. \textbf{144} (2004), 221--291.
\href{https://doi.org/10.1007/BF02916714}{doi:10.1007/BF02916714}.

\bibitem{SabbahFourier}
C.~Sabbah,
\emph{Fourier transformation and Stokes structures},
lecture notes, version of February 13, 2020,
\url{https://perso.pages.math.cnrs.fr/users/claude.sabbah/livres/sabbah_chicago1205.pdf}.

\end{thebibliography}
\end{document}
